\documentclass[a4paper, french, twoside, twocolumn, 11 pt]{article}

% Language setting
% Replace `english' with e.g. `french' to change the document language
\usepackage[french]{babel}

% Set page size and margins
% Replace `letterpaper' with `a4paper' for UK/EU standard size
\usepackage[a4paper,top=2cm,bottom=2cm,left=1.4cm,right=2cm,marginparwidth=1.75cm]{geometry}

% Useful packages
\usepackage{graphicx}
\usepackage{siunitx}
\usepackage[colorlinks=true, allcolors=blue]{hyperref}
\usepackage{pgfplots}
%\usepackage{tkz-euclide}
\usepackage[T1]{fontenc}
\usepackage[utf8]{inputenc}
%\usepackage[european, straightvoltages, RPvoltages]{circuitikz}
%\usetikzlibrary{babel}
\usepackage{indentfirst}
%\usepackage{esint}
\usepackage{fancyhdr}
%\\
\usepackage{amsfonts, amsmath, amssymb}
\usepackage{subcaption}
%\usepackage[export]{adjustbox}
\usepackage{lipsum}
%\usepackage{mathtools}
%\usepackage{tkz-tab}
%\usepackage{longtable}
%\usepackage[cyr]{aeguill}
%\usepackage{xspace}
%\usepackage{titling}
%\usepackage[absolute,overlay]{textpos}
%\usepackage{etoc}
\usepackage{enumitem}
\usepackage{multirow}
\usepackage{nicefrac}
\usepackage{sectsty}
\usepackage{titlesec}
\usepackage[autostyle=true, french=guillemets]{csquotes}
\usepackage{appendix}
\usepackage{mathrsfs}
\usepackage{mathabx}
\usepackage{bm}
%\pgfplotsset{compat=newest}

%Mathematical operators
\DeclareMathOperator{\e}{e}
\newcommand{\deriv}{\mathrm{d}}

%Bibliography
\usepackage[
backend=biber,
style=numeric-comp,
sorting=nty
]{biblatex}
\addbibresource{Reference_stage.bib}

\usepackage{nameref}
\makeatletter
\newcommand*{\currentname}{\@currentlabelname}
\makeatother

%Style
\pagestyle{fancy}
\fancyhf{}
\fancyhead[C]{\textsc{Rapport de stage}}
\fancyhead[R]{}
\fancyfoot[RO, LE]{\emph{\thepage}}

\fancypagestyle{1erepage}{
\renewcommand{\headrulewidth}{0pt}
\renewcommand{\footrulewidth}{0pt}
\fancyhead[RO, LE]{\footnotesize{Compilé avec pdf\LaTeX}}
\fancyhead[C]{}
\fancyfoot[RE, LO]{}
\fancyfoot[RO, LE]{}
}

\def\changemargin#1#2{\list{}{\rightmargin#2\leftmargin#1}\item[]}
\let\endchangemargin=\endlist

\titlespacing*{\section}
{0pt}{2.5ex plus 1ex minus .0ex}{0.0ex plus .0ex}
\allsectionsfont{\centering}

%\titleformat{\subsection}{\normalfont\normalsize\normalseries}
\titlespacing*{\subsection}
{0pt}{1.5ex plus 0ex minus .0ex}{0.0ex plus .0ex}

\begin{document}
\newpage
\thispagestyle{1erepage}
\title{\textbf{\Large{Stage Magistère}\\ \LARGE{Modélisation d'un gaz chaud au sein d'un halo galactique}\\[0.5em]}}
\author{\normalsize\textbf{Tibor \textsc{Donin de Rosière}}\\\vspace{-0.5em}
    \footnotesize\emph{Étudiant en L3 Magistère de Physique}\\\vspace{-0.5em}
    \footnotesize\emph{Université Grenoble Alpes, Laboratoire de Physique Subatomique et de Cosmologie, F-38026 Grenoble, France}}
    \date{}

\twocolumn[{
\vspace{-3em}
\maketitle
\begin{abstract}
    %\vspace{-1em}
        %\begin{spacing}{0.9}
            \emph
            {Dans les halos galactiques, nous pouvons trouver la présence d'un noyau galactique, d'une couronne de gaz et plasma chaud et d'un halo de matière noire. Ce gaz chaud est essentiellement composé de noyau d'hydrogène, d'hélium et d'électrons. Ces derniers peuvent, via l'effet Compton inverse, interagir avec les photons du fond diffus cosmologique (CMB) et provoquer des fluctuations de température dans les observations de ce rayonnement. Ces fluctuations sont liées à la distribution de pression et à la densité du gaz d'électrons. Par un modèle théorique de la fluctuation de température, il est possible d'en déduire l'activité du noyau galactique dans le halo ainsi que l'impact gravitationnel par la matière noire sur la répartition du gaz.}
        %\end{spacing}
\end{abstract}
\vspace{3em}
}]
\thispagestyle{1erepage}
\section{Introduction}
La formation d'halos galactiques est un phénomène se déroulant à différentes échelles. Les champs gravitationnelles à grandes échelles contraignent les trajectoires de la matière, formant de riches structures composées de matière noire et ordinaire, plus communément appelé baryonique, tel que des filaments ou des galaxies. On définit donc le halo comme une composante sphérique d'une galaxie qui s'étend au-delà de la partie principale visible de la galaxie. Il est constitué d'un halo stellaire, d'une couronne galactique composé de gaz chaud et d'un halo de matière noire. De plus, les noyaux galactiques actifs (AGN) aux centres des galaxies massives et les supernovaes libèrent de grandes quantités d'énergie sous la forme de jets et de vents galactiques très rapides pouvant impacter la distribution de baryons initialement présente. Du fait que ces jets et vents rapides poussent le gaz chaud diffus hors de l'halo, perturbant sa température et sa densité, cela impacte gravitationnellement la distribution de matière noire. Il y a donc un effet de rétroaction corrélé avec cet expulsion du gaz chaud sur la distribution de l'ensemble de la matière \cite{pandey2024godmaxmodelinggasthermodynamics}.\\
Par le biais de la diffusion Compton, le gaz ionisé présent autour des galaxies et des amas laisse plusieurs empreintes distinctes sur le rayonnement du fond diffus cosmologique (CMB). Les deux effets principaux sont les décalages Doppler des photons du CMB dus à la dispersion des vitesses du gaz,l'effet thermique de Sunyaev-Zeldovich (tSZ), et ceux dus au mouvement global du gaz, c'est-à-dire l'effet cinématique de Sunyaev-Zeldovich (kSZ).
Ces deux effets sont utilisés comme observables \cite{Birkinshaw_1986} pour rendre compte de la fluctuation de la densité du gaz chaud et notamment de ses électrons qui le compose. Le kSZ possède moins de mesures observées ce qui m'a mené à poser la problématique de l'impact des électrons libres dans le gaz chaud diffus sur le rayonnement du CMB. Pour cela, je modélise l'effet kSZ en traçant l'évolution de sa fluctuations de température $T_\text{kSZ}$ en fonction de l'écartement angulaire $\theta$ par rapport au centre du halo, etle comparer qualitativement avec des ajustements déjà existants. Pour réaliser ce modèle, je considère mon système comme étant un halo galactique de forme sphérique avec un rayon $R\simeq0,1\ \unit{Mpc}$, et de masse $M=10^{13}\ \unit{M_\Sun}$. La distribution du gaz utilisé est celui du modèle dit de baryonification de \cite{schneider2025baryonificationalternativehydrodynamicalsimulations} car il possède plusieurs paramètres ajustables notamment les valeurs de pentes dans différentes régions de la couronne du gaz. J'utilise les hypothèses du modèle cosmologique $\Lambda$CDM qui est actuellement la description de l'univers la plus en vogue, d'après les mesures de \cite{Planck_2020}. 

La structure de ce rapport est la suivante : dans la section \ref{Section2}, je décris le modèle théorique d'halo galactique et de gaz chaud avec les hypothèses admises et leurs éventuelles justifications, et les calculs des observables. Je décris également comment ces observables ont été implantées dans le modèle en question. En section \ref{Section3}, je présente les résultats de ces modèles et les discute, et dans la Section \ref{Section4}, je conclus.
\section{Méthodes}
\label{Section2}
L'effet Sunyaev-Zeldovich peut être divisé en deux types différents, thermique et cinématique.
Ces deux signaux peuvent être modélisés en termes de fluctuations de température comme décrit ci-dessous.
\subsection{Sunyaev-Zeldovich thermique} 
L'effet Sunyaev-Zeldovich thermique (tSZ) est une distorsion spectrale du CMB produite par la diffusion des photons du CMB sur des électrons chauds du milieu intra-amas (ICM) via l'effet Compton inverse. Ce phénomène fut prédit en 1972 par Rashid Sunyaev et Yakov Zel'dovich. Lorsque les photons du CMB traversent un amas de galaxies contenant du gaz ionisé à haute température ($T_\text{e}\simeq10^7-10^8\ \unit{K}$), ils interagissent avec les électrons énergétiques par diffusion Compton inverse.
Les fluctuations de température dues à la dispersion des vitesses du gaz sont données par
\begin{equation}
\frac{\Delta T_\text{tSZ}}{T_\text{CMB}}=f(x)\, y
\end{equation}
où $f(x) = x\coth\!\frac{x}{2}-4$ est la dépendance vis-à-vis de la fréquence tel que $x=\frac{h\nu}{k_BT}$, sans prendre en compte les corrections relativistes. $T_\text{CMB}$ est la température du CMB, $k_B$ la constante de Boltzmann et $h$ la constante de Planck. Pour plus de détail sur le paramètre Compton $y$, se référer à l'appendice \ref{AppendixA}.

\subsection{Sunyaev-Zeldovich cinématique}
L'effet Sunyaev-Zeldovich cinématique (kSZ) est quand à lui, un effet secondaire où les photons du CMB interagissent avec des électrons de hautes énergies dû au mouvement du bulbe galactique. Il est proportionnel à la densité numérique d'électron $n_\text{e}$ et à la vitesse particulière $v_p$ du gaz.
Cette diffusion inverse modifie la température du CMB selon, d'après \cite{Sunyaev_1972}, la formule suivante
\begin{equation}
\label{eqDeltakSZ}
\frac{\Delta T_\text{kSZ}}{T_\text{CMB}}=\frac{\sigma_\text{T}}{c}\int_\text{los} \deriv l\,e^{-\tau} n_\text{e} \; v_\text{p}
\end{equation}
où $\sigma_\text{T}$ est la section efficace de Thomson. Je suppose que le gaz se déplace de manière cohérente avec la vitesse globale du halo (voir réf. \cite{Hadzhiyska2023}), ce qui permet d'éliminer la vitesse particulière du gaz $v_p$ de l'intégrale dans la ligne de visée. Cette simplification est couramment utilisée dans les analyses par empilement, où les vitesses individuelles ne sont pas connues avec précision. Au lieu de cela, le signal attendu est modélisé à l'aide de la vitesse radiale quadratique moyenne (RMS) de l'échantillon de halos, que je fixe à $v_r = 1,06 \times 10^{-3}c$, conformément à la référence \cite{Amodeo2021}. Le calcul est plus détaillé dans l'appendice \ref{AppendixC}.\\
La densité numérique d'électrons dans l'équation \ref{eqDeltakSZ} peut être calculée à partir de
\begin{equation}
n_\mathrm{e} = \frac{X_\mathrm{H} + 1}{2 m_\mathrm{amu}} \rho_\mathrm{gas}
\end{equation}
où $X_\mathrm{H}$ est la fraction massique d'hydrogène dans l'univers et $m_\mathrm{amu}$ est l'unité atomique de masse. La densité massique du gaz $\rho_\mathrm{gas}$ est obtenu à partir du modèle de baryonification de \cite{schneider2025baryonificationalternativehydrodynamicalsimulations}.
\begin{equation}
\begin{split}
& \rho_\text{gas}(r) = \rho_\text{gas,0}\left[1+\left(\frac{r}{\theta_\text{c}r_{500c}}\right)^{\alpha}\right]^{-\beta(M_{500})/\alpha} \\
& \left[1+\left(\frac{r}{\varepsilon r_{500c}}\right)^{\gamma}\right]^{\frac{\delta}{\gamma}}
\end{split}
\end{equation}
Ici, $\theta_\text{c}$ est le paramètre du rayon du noyau, $\epsilon$ le paramètre du rayon tronqué avec $\varepsilon=\varepsilon_0+\varepsilon_1\nu$ où $\nu=\frac{\delta_c}{\sigma(M_{500})}$ est le peak height du halo. Par simplification du modèle, j'utilise les valeurs admises dans \cite{kovac2025baryonificationiiconstrainingfeedback} avec $\theta_\text{c}=0,3$; $\varepsilon_0 = 4$; $\varepsilon_1=0,5$ et $\delta_c=1,686$. Ensuite, $\beta(M_{500})$ détermine la pente du profil tel que
\begin{equation}
\beta(M_{500})=\frac{3(M_{500}/M_\text{c})^{\mu}}{1+(M_{500}/M_\text{c})^{\mu}}
\end{equation}
où $M_\text{c}$ et $\mu$ sont des paramètres libres contrôlant la masse pivot et le paramètre de transition, et les paramètre $\alpha$, $\delta$ et $\gamma$ définissent les pentes intérieure et extérieure et la région de transition. Les valeurs de $\alpha$ et $\gamma$ sont fixés à 0,3081 et à 1,0510 conformément aux meilleurs ajustements des simulations de \cite{Arnaud_2010} avec les observations.\\
Pour modéliser des mesures, on considère le profil kSZ comme un signal où l'on applique une convolution sur l'Eq. \ref{eqDeltakSZ} comme utilisé dans \cite{Schaan_2021}. De plus, afin d'atténuer le bruit dû aux variations à grande échelle du rayonnement CMB, un filtre de photométrie à ouverture compensée a été appliqué au modèle. Ce filtre est une fonction de pondération $W_{\theta_\text{d}}(\theta)$ prenant la forme
\[
W_{\theta_\text{d}}(\theta) =
\begin{cases}
+1 & \text{si } \theta < \theta_\text{d} \\
-1 & \text{si } \theta_\text{d} \leq \theta < \sqrt{2}\, \theta_\text{d} \\
0 & \text{sinon}
\end{cases}
\]
où $\theta_\text{d}$ est le rayon d'ouverture. Le modèle s'accorde aux points des données issus des halos \enquote{constant mass} CMASS empilés à 150 GHz de \cite{Schaan_2021}. À cette fréquence, on a $\theta_\text{d}=1,3\ \unit{arcmin}$\\
On calcule ensuite, la température kSZ en utilisant
\begin{equation}
    \mathcal{T}(\theta_\text{d})=\int \deriv^2{\bm{\theta}}\,\delta T({\bm{\theta}})\,W_{\theta_\text{d}}({\bm{\theta}})\: ,
\end{equation}
Le résultat est généralement exprimé en $\mu$K·arcmin$^2$, ce qui reflète une intégration sur l'angle solide qui est la convention adoptée dans la plupart des analyses des observations\cite{bigwood2025kineticsunyaevzeldovicheffect, Hadzhiyska2023}. Le code de la modélisation, générant les figures, est disponible sur le dépôt github \href{https://github.com/AdrikForgefeu/Stage-Magistere-LPSC.git}{Stage-Magistere-LPSC}.

\section{Résultats}
\label{Section3}
L'effet kSZ a été modélisé pour un halo galactique à un redshift $z=0,5$. La densité volumique au centre du halo $\rho_\text{gas,0}$ est estimé à $10^{15}\ \unit{M_\Sun.Mpc^{-3}}$ depuis une valeur de masse du gaz de $m_\text{gas}=10^9\ \unit{M_\Sun}$ \cite{kovac2025baryonificationiiconstrainingfeedback} sur un noyau galactique de rayon $R=10\ \unit{kpc}$. Le tSZ étant un effet ayant plus de résultats observationels et est étudié plus en profondeur, il possède plus de modèle existant. Les résultats montrés ici, ne se concentrent donc que sur le kSZ. Ils sont résumés par la figure \ref{figkSZ}, dans l'objectif de pouvoir vérifier qualitativement la validité de l'allure avec d'autres études. 
Ce modèle possède deux paramètres significatifs, la masse pivot $M_c$ et la pente extérieure $\delta$ du profil. $\mu$ n'a aucune influence sur le kSZ. La figure \ref{figkSZ} présente l'évolution de la température sur surface angulaire $T_\text{kSZ}$ en fonction de l'écartement angulaire $\theta$ pour différentes masses pivot. La masse pivot définit l'échelle de masse caractéristique à laquelle se produit la transition de $\beta(M_{500})$. Ce paramètre influence la rapidité avec laquelle le profil de densité du gaz s'accentue à mesure que la masse du halo augmente. Pour une masse $M_{500c}$ fixe, $\beta$ tend vers 0 pour de hautes masses pivots. Le terme correspondant à la partie interne de la couronne galactique ne joue pas sur l'intégration de $n_\text{e}$.\\
À petit $\delta$, on observe que l'asymptote de la pente de $T_\text{kSZ}$ augmente pour de grand écartements angulaires. En effet, celle-ci joue sur la rapidité de la décroissance de $\Delta T_\text{kSZ}$.\\
En comparant l'allure graphique avec \cite{bigwood2025kineticsunyaevzeldovicheffect, kovac2025baryonificationiiconstrainingfeedback, Oppenheimer_2025}, j'obtiens les mêmes ordres de grandeurs pour la convergence qui sont aux alentours de $10\ \unit{\mu K.arcmin^2}$.
\begin{figure}[t]
\centering
\includegraphics[width=0.405\textwidth]{temp_ksz_profile.png}
\caption{\footnotesize{Evolution de l'effet kSZ avec l'écartement angulaire en fonctions des paramètres $M_c$ et $\delta$ pour un redshift de $z=0,5$. Les deux graphiques montrent comment la variation d'un seul paramètre du modèle influe sur la valeur observable prédite, tandis que l'autre paramètre est maintenu constant à la valeur médiane de sa rangée de valeurs $M_\text{c}=10^{13}\ \unit{M_\Sun}$, $\delta=5$ et $\mu=1$. Le paramètre $\mu$ ne possède pas un impact significatif sur l'allure du profil.}}
\label{figkSZ}
\end{figure}

\section{Conclusion}
\label{Section4}
Lorsque le rapport $\nicefrac{M_{500}}{M_\text{c}}$ diminue, cela signifie que la densité du gaz est faible au centre du halo. Cette absence de gaz peut être dû à une rétroaction passé, ayant éjecté le gaz présent plus vers la périphérie. De plus, la variation $\Delta T_\text{kSZ}$ est moins importante. La présence de noyau actif  possède donc impact indirect sur l'intensité de la fluctuation de température aux abords du halo.\\ 
À partir d'une certaine pente $\delta$ on observe physiquement que au dessus de cette valeur, les vents stellaires libérées par les AGNs sont suffisamment puissants pour avoir un effet sur la répartition du gaz.
Une petite valeur de $\delta$ entraîne un profil $n_\text{e}$ plus plat qui s'étend jusqu'à de grands rayons, ce qui redistribue davantage de gaz vers la périphérie et réduit la densité centrale.
Cela entraîne une atténuation du signal kSZ au centre et une diminution des fractions de gaz à l'intérieur de $R_{500}$, car la majeure partie de la masse se trouve au-delà de ce rayon. Dans ce cas, la distribution du gaz suit un profil de densité volumique Navarro-Frenk-White par interaction gravitationnelle avec la matière noire.\\
Cette étude a mit en évidence le comportement du gaz d'électrons au sein de la couronne galactique, et permit l'identification d'une activité hydrodynamique au sein du noyau. Et tout cela par la modélisation des fluctuations de la température du CMB.

\section*{Remerciements}
En premier lieu, je souhaite remercier Mme. Marine Kuna qui m'a gentiment mis en contact avec Mr. Johan Comparat mon tuteur. Je remercie ce dernier pour la proposition du sujet passionnant et pédagogique m'ayant permis d'apprendre à être rigoureux et à créer mes propres outils de travail lors d'une démarche méthodologique de recherche. Et enfin, je tiens à remercier l'ensemble de l'équipe Cosmologie observationnelle pour m'avoir accueilli au sein de leurs locaux.

\medskip
\printbibliography[
title={Références}
]
\newpage
\appendix
\appendixpage
\section{Paramètre Compton-$y$}
\label{AppendixA}
L'amplitude résultante de cette distorsion spectrale est le paramètre adimensionnel Compton-$y$ \cite{Sunyaev_1972} mesuré selon l'écartement angulaire $\theta$, à une distance angulaire d'un redshift $z$, $d_A(z)$ est
\begin{equation}
\begin{split}
& y = \int \sigma_\text{T}\, n_\text{e}\, \frac{k_B T_\text{e}}{m_\text{e} c^2}\, \deriv l \\
& = \frac{\sigma_\text{T}}{m_\text{e} c^2}\int P_\text{e}\left(\sqrt{l^2+d_A(z)^2|\theta|^2}\right)\, \deriv l
\end{split}
\end{equation}
où $P_\text{e}=n_\text{e}k_BT_\text{e}$ la pression électronique et $\deriv l$ est la distance physique de la ligne de visée. La distance $r$ par rapport au noyau galactique peut être décomposé comme $r=\sqrt{l^2+d_A(z)^2|\theta|^2}$.
La pression thermique de l'ICM suit une courbe de profil Navarro-Frenk-White généralisée quasi universelle (\cite{Nagai_2007}), calibrée sur l'échantillon REXCESS par \cite{Arnaud_2010}. Pour plus de précision concernant le profil électronique $P_\text{e}$, voir appendice \ref{AppendixB}.
\begin{figure}[h]
\centering
\includegraphics[width=0.45\textwidth]{tsz_y_profile.png}
\caption{\footnotesize{Evolution du Compton-$y$ en fonction de l'écartement angulaire pour une masse $M_{500}=10^{14}\ \unit{M_\Sun}$ à un redshift $z=0,5$. $y$ traduit la décroissance de la pression électronique le long de la ligne de visée ($-5R_{500}<l<5R_{500}$)}}
\end{figure}

\section{Profil de pression thermique}
\label{AppendixB}
La pression de l'ICM obtenu par \cite{Arnaud_2010} est de la forme suivante
\begin{equation}
P(r) = P_{500}\; \text{p}(x), \qquad x = r/R_{500},
\end{equation}
où $R_{500}$ représente le rayon du halo à 500 fois la densité critique de l'univers $\rho_0$ et avec
\begin{equation}
\text{p}(x) = \frac{P_0}{(c_{500}x)^{\gamma}\,\big[1+(c_{500}x)^{\alpha}\big]^{(\beta-\gamma)/\alpha}}
\end{equation}
avec les paramètres d'ajustements optimaux $[P_0;c_{500};\gamma;\alpha;\beta] = [8,403\,h_{70}^{-3/2};\,1,177;\,0,3081;\,1,0510;\,5,4905]$ \cite{Arnaud_2010}. La pression caractéristique est
\begin{equation}
\begin{split}
& P_{500} =  1,65\times10^{-3}\,E(z)^{8/3}\times \\
& \left(\frac{M_{500}}{3\times10^{14}\,h_{70}^{-1}M_\odot}\right)^{2/3} h_{70}^{2}\;\;{\rm keV\,cm^{-3}}
\end{split}
\end{equation}
$E(z)=\sqrt{\Omega_\text{m}(1+z)^3+(1-\Omega_\text{m})}$ représente le paramètre de Hubble normalisé avec $\Omega_\text{m}$ le paramètre de densité pour la matière.
Le tSZ est corrélé à une pression $P_e = (\mu/\mu_e)\,P$ avec les masses moléculaires moyennes $\mu\simeq0,59$ et $\mu_e\simeq1,14$ pour un plasma entièrement ionisé d'hydrogène et d'hélium.
\begin{figure}[h]
\centering
\includegraphics[width=0.45\textwidth]{tsz_Pe_profile.png}
\caption{\footnotesize{Evolution de la pression électronique $P_\text{e}$ en fonction de la distance au centre du halo $r$ d'après le profil de \cite{Arnaud_2010} à $M_{500}=10^{14}\ \unit{M_\Sun}$ et $z=0,5$ }}
\end{figure}

\section{Profondeur optique}
\label{AppendixC}
La profondeur optique $\tau$ le long de la ligne de visée due à la diffusion de Thompson entre l'observateur et le redshift $z$ est définit comme suit :
\begin{equation}
\tau(\theta) =  \sigma_\mathrm{T} \int_{\rm los} n_{\mathrm{e}}\left(\sqrt{l^2\:+\:d_A(z)^2\theta^2}\right)\deriv l\;,
\end{equation}
Elle mesure la quantité absorbée de lumière du CMB par le gaz chaud d'électrons par rapport au CMB incident. La profondeur optique moyenne dans l'intervalle de valeur de redshift étudié (0,4 < $z$ < 0,7) est inférieure à 0,01 (voir \cite{Adam_2016}), j'effectue donc l'approximation suivante $e^{-\tau}\approx1$. On peut simplifier plus profondément l'équation \ref{eqDeltakSZ} en
\begin{equation}
\frac{\Delta T_{\mathrm{kSZ}}}{T_{\mathrm{CMB}}}=\tau_{\mathrm{gal}}\left(\frac{v_\mathrm{r}}{c}\right).
\end{equation}
où $\tau_\text{gal}$ est la profondeur optique après intégration sur une ligne de visée correspondant à la taille d'une galaxie. En effet, le SZ se prononce essentiellement au centre de l'halo et par conséquent au niveau de la galaxie le contenant.
\begin{figure}[h]
\centering
\includegraphics[width=0.45\textwidth]{ksz_tau_profile.png}
\caption{\footnotesize{Evolution de la profondeur optique $\tau$  selon l'écartement angulaire sur un halo galactique de $M_{500}=10^{13}\ \unit{M_\Sun}$ à un redshift $z=0,5$. $\tau$ est proportionnelle à la densité numérique d'électrons $n_\text{e}$, ce qui implique une présence du gaz chaud d'électrons plus importante au niveau du noyau galactique.}}
\end{figure}
\end{document}